{\bfseries Comparación con el PageRank clásico. 
Para el mismo grafo ciclo de 3 vértices, calcule el PageRank clásico con factor de amortiguamiento $\alpha=0.85$. 
Recuerde que la matriz de transición $M$ se construye a partir de los grados, y la matriz de PageRank es: 
\begin{equation*}
P=\alpha M+(1-\alpha)\frac{1}{3}\mathbf{1}\mathbf{1}^{T}
\end{equation*}}

\begin{enumerate}
    \item[a)] {\bfseries Escriba explícitamente la matriz $P$.}
    \item[b)] {\bfseries Encuentre el vector estacionario $\pi$ (único) que satisface $P\pi=\pi$ con $\sum \pi_{i}=1$.}
    \item[c)] {\bfseries Compare los valores de $\pi_{i}$ con los $QR_{i}$ obtenidos en el ejercicio anterior. ¿Hay diferencias significativas? ¿Cuál de los dos métodos asigna mayor importancia al vértice central?}
\end{enumerate}

\textbf{a) Matriz $P$ explícita:}
Dado que el grafo es un ciclo, cada vértice tiene grado $d=2$. La matriz de transición clásica estocástica por columnas $M$, definida como $M_{ij} = A_{ij}/d_j$, es:
\begin{equation*}
M = \begin{pmatrix} 0 & 1/2 & 1/2 \\ 1/2 & 0 & 1/2 \\ 1/2 & 1/2 & 0 \end{pmatrix}
\end{equation*}
La matriz de teletransportación es $\frac{1}{3}\mathbf{1}\mathbf{1}^T = \begin{pmatrix} 1/3 & 1/3 & 1/3 \\ 1/3 & 1/3 & 1/3 \\ 1/3 & 1/3 & 1/3 \end{pmatrix}$.
Sustituyendo el factor de amortiguamiento $\alpha = 0.85$:
\begin{equation*}
P = 0.85 \begin{pmatrix} 0 & 0.5 & 0.5 \\ 0.5 & 0 & 0.5 \\ 0.5 & 0.5 & 0 \end{pmatrix} + 0.15 \begin{pmatrix} 1/3 & 1/3 & 1/3 \\ 1/3 & 1/3 & 1/3 \\ 1/3 & 1/3 & 1/3 \end{pmatrix}
\end{equation*}
\begin{equation*}
P = \begin{pmatrix} 0 & 0.425 & 0.425 \\ 0.425 & 0 & 0.425 \\ 0.425 & 0.425 & 0 \end{pmatrix} + \begin{pmatrix} 0.05 & 0.05 & 0.05 \\ 0.05 & 0.05 & 0.05 \\ 0.05 & 0.05 & 0.05 \end{pmatrix} = \begin{pmatrix} 0.05 & 0.475 & 0.475 \\ 0.475 & 0.05 & 0.475 \\ 0.475 & 0.475 & 0.05 \end{pmatrix}
\end{equation*}

\vspace{0.5cm}
\textbf{b) Vector estacionario $\pi$:}
Como el grafo es regular, la matriz de adyacencia y la de transición son simétricas. Esto hace que la matriz $P$ sea doblemente estocástica (la suma de los elementos de cualquier fila o columna es igual a 1). Para toda matriz doblemente estocástica irreducible de dimensión $N \times N$, el único vector estacionario $\pi$ es la distribución uniforme uniforme:
\begin{equation*}
\pi = \begin{pmatrix} 1/3 \\ 1/3 \\ 1/3 \end{pmatrix}
\end{equation*}
De este modo se satisface inmediatamente que $P\pi = \pi$ y $\sum_{i} \pi_i = 1/3 + 1/3 + 1/3 = 1$.

\vspace{0.5cm}
\textbf{c) Comparación e interpretación:}
\begin{itemize}
    \item \textbf{Valores:} En el PageRank clásico tenemos $\pi_0 = \pi_1 = \pi_2 = 1/3 \approx 0.333$. Por el contrario, en el PageRank cuántico obtuvimos $QR_0 = 5/9 \approx 0.555$ y $QR_1 = QR_2 = 2/9 \approx 0.222$.
    \item \textbf{Diferencias significativas:} Son evidentes y se deben a la topología y las condiciones iniciales. El PageRank clásico refleja la estructura del grafo: como todos los nodos tienen la misma conectividad (son topológicamente equivalentes), todos reciben la misma importancia. El PageRank cuántico, debido a su dinámica de evolución unitaria (sin decoherencia y reversible), mantiene una fuerte dependencia de la condición inicial $\ket{\psi(0)} = \ket{0}$, concentrando fuertemente la probabilidad allí debido a la interferencia cuántica.
    \item \textbf{Importancia al vértice central/inicial:} Como lo insinúa el problema, en un ciclo de 3 vértices no existe un "vértice central" estructuralmente hablando; todos son periféricos o todos son centrales. Sin embargo, refiriéndonos al vértice 0 como el centro del análisis por ser el origen, el \textbf{PageRank cuántico} es el método que le asigna una importancia significativamente mayor ($55.5\%$ frente al $33.3\%$ clásico).
\end{itemize}
